% =============================================================
%  Kapitel 7 – Traceability: Anforderung → Artefakt
% =============================================================
\section{Traceability: Anforderung \texorpdfstring{$\to$}{->} Artefakt}
\label{sec:traceability}

Die folgende Tabelle ordnet jede Pflichtanforderung der Aufgabenstellung
einem konkreten Artefakt in diesem Dokument zu.

\begin{center}
\renewcommand{\arraystretch}{1.4}
\begin{longtable}{p{6.5cm} p{4.5cm} p{2.5cm}}
  \toprule
  \textbf{Anforderung (Aufgabenstellung)} &
  \textbf{Artefakt / Diagramm} &
  \textbf{Abschnitt} \\
  \midrule
  \endhead
  % --- Architektur ---
  Architektur-Übersicht (Context Diagram)
    & Abb.~\ref{fig:context} \newline 01\_Context.png
    & \S\,\ref{sec:systemkontext} \\
  Subsystem-Zerlegung + Verantwortlichkeiten
    & Abb.~\ref{fig:component_overview} \newline 02\_Component\_Overview.png
    & \S\,\ref{sec:component_overview} \\
  Class Diagrams (Jetson)
    & Abb.~\ref{fig:jetson_classes} \newline 02\_Jetson\_ClassDiagram.png
    & \S\,\ref{sec:klassendiagramme} \\
  Class Diagrams (Master)
    & Abb.~\ref{fig:master_classes} \newline 02\_Master\_ClassDiagram.png
    & \S\,\ref{sec:klassendiagramme} \\
  Class Diagrams (Slave)
    & Abb.~\ref{fig:slave_classes} \newline 02\_Slave\_ClassDiagram.png
    & \S\,\ref{sec:klassendiagramme} \\
  % --- Deployment ---
  Hardware/Software-Mapping (Deployment)
    & Abb.~\ref{fig:deployment} \newline 03\_Deployment.png
    & \S\,\ref{sec:deployment} \\
  % --- Global Software Control ---
  Global Software Control (Sequenz Safety-Zone)
    & Abb.~\ref{fig:seq_safetyzone} \newline 04\_Sequence\_SafetyZone.png
    & \S\,\ref{sec:seq_safetyzone} \\
  % --- Boundary Conditions ---
  Boundary Condition: Startup
    & Abb.~\ref{fig:seq_startup} \newline 05\_Sequence\_Startup.png
    & \S\,\ref{sec:seq_startup} \\
  Boundary Condition: Shutdown
    & Abb.~\ref{fig:seq_shutdown} \newline 05\_Sequence\_Shutdown.png
    & \S\,\ref{sec:seq_shutdown} \\
  Boundary Condition: Error
    & Abb.~\ref{fig:seq_error} \newline 05\_Sequence\_Error.png
    & \S\,\ref{sec:seq_error} \\
  Boundary Condition: Release
    & Abb.~\ref{fig:seq_release} \newline 05\_Sequence\_Release.png
    & \S\,\ref{sec:seq_release} \\
  Boundary Condition: Release Rejected
    & Abb.~\ref{fig:seq_releaserejected} \newline 05\_Sequence\_ReleaseRejected.png
    & \S\,\ref{sec:seq_releaserejected} \\
  % --- Detailed Design ---
  Detailed Design ES: Klasse + private Methoden
    & Abb.~\ref{fig:estop_class} \newline 06\_EStop\_Class\_Detailed.png + Tabellen
    & \S\,\ref{sec:estop_klassen} \\
  Detailed Design ES: Internes Sequenzdiagramm
    & Abb.~\ref{fig:estop_seq} \newline 06\_EStop\_Internal\_Sequence.png
    & \S\,\ref{sec:estop_seq} \\
  Detailed Design ES: State Machine
    & Abb.~\ref{fig:estop_statemachine} \newline 06\_EStop\_StateMachine.png
    & \S\,\ref{sec:estop_statemachine} \\
  % --- Entscheidungen ---
  Sprachauswahl (C vs.\ C++) + Rationale
    & Text + Tabelle
    & \S\,\ref{sec:sprachauswahl} \\
  Architektur-Alternativen + Rationale
    & Text
    & \S\,\ref{sec:entsch_verteilung} \\
  % --- Offene Punkte ---
  Activity-Diagramme pro private Methode
    & \todo{Ausstehend}
    & \S\,\ref{sec:naechste_schritte} \\
  \bottomrule
\end{longtable}
\end{center}
